\section{Критерий полноты нормированного пространства.}
\begin{theorem}
    Пусть дано нормированное пространство $(X, \|\cdot\|)$.
    Следующие условия эквивалентны:
    \begin{enumerate}
        \item $X$ -- банахово.
        \item Всякий абсолютно сходящийся в $X$ ряд сходится.
    \end{enumerate}
\end{theorem}
\begin{proof}
    Пусть $X$ -- банахово.
    Тогда если $\sum_{i = 1}^{\infty} x_i$ -- абсолютно сходящийся ряд, то есть
    \[
        \sum_{i = 1}^{\infty} \|x_i\| < +\infty,
    \]
    то последовательность норм -- фундаментальна и, следовательно последовательность частичных сумм ряда также фундаментальна:
    \[
        \|S_n - S_m\| \leq \sum_{i = n + 1}^{m} \|x_i\|.
    \]
    В силу полноты $X$ она сходится к некоторому $x \in X$.

    Пусть всякий абсолютно сходящийся ряд сходится.
    Рассмотрим фундаментальную последовательность $\{x_n\}$.
    Выберем из неё подпоследовательность $\{x_{n_k}\}$ такую, что
    \[
        \|x_{n_k} - x_m\| \leq 2^{-k}, \ \forall m \geq n_k.
    \]
    Пусть теперь $y_0 = x_{n_1}$, а $y_k = x_{n_{k + 1}} - x_{n_k}, k \geq 1$.
    Заметим, что ряд $\sum_{m = 0}^{\infty} y_m$ -- абсолютно сходится:
    \[
        \sum_{m = 0}^{\infty} \|y_m\| \leq \|x_{n_1}\| + \sum_{m = 1}^{\infty} 2^{-m} < +\infty.
    \]
    Тогда он и просто сходится, то есть $\exists x \in X\colon$
    \[
        x = \sum_{m = 0}^{\infty} y_0.
    \]
    Заметим, что частичные суммы ряда имеют вид $S_k = x_{n_{k + 1}}$.
    А значит подпоследовательность $x_{n_k}$ сходится к некоторому элементу $X$, и, в силу фундаментальности, и вся последовательность сходится к этому элементу.
    Следовательно, $X$ -- банахово.
\end{proof}
\begin{theorem}[Полнота $L^\infty$.]
    Пусть $(X, \mathcal{M}, \mu)$ -- пространство с мерой.
    Тогда пространство $L^\infty(X, \mathcal{M}, \mu)$ является банаховым.
\end{theorem}
\begin{proof}
    Пусть $\{f_n\} \subset L^{\infty}$ -- фундаментальная последовательность.
    Заметим, что для почти всех $x \in X$ последовательность $\{f_n(x)\}$ -- фундаментальна.
    Тогда у нас почти для всех $x$ корректно определена функция $f(x) = \lim f_n(x)$.
    Заметим, что тогда в силу фундаментальности $f_n\colon$
    \[
        \forall \epsilon > 0 \ \exists N \in \N \ \forall m, n \geq N \ \exists E_{m, n}\colon \mu(E_{m, n}) = 0 \hookrightarrow |f_n(x) - f_m(x)| \leq \epsilon.
    \]
    Пусть $E_n = \bigcup_{m \geq N} E_{m, n}$.
    Оно является множеством меры нуль.
    Вне его корректен предельный переход:
    \[
        |f_n(x) - f(x)| \leq \epsilon, \ \forall x \notin E_n.
    \]
    Из чего и следует $\|f_n - f\|_\infty$ -- что и означает сходимость к $f$ по $\esssup$-норме (и из чего очевидным образом следует и существенная ограниченность $f$).
\end{proof}
\section{Замкнутость конечномерных подпространств.}
\begin{proposition}
    Пусть $(L, \|\cdot\|)$ -- конечномерное нормированное пространство.
    Тогда $L$ -- банахово.
\end{proposition}
\begin{proof}
    Выберем некоторый базис $e_1, \ldots, e_n$.
    Всякая точка $x$ из $L$ раскладывается по нему единственным образом:
    \[
        x = \sum_{i = 1}^{n} x_i e_i.
    \]
    Определим $\|x\|_e$ как
    \[
        \|x\|_e = \sum_{i = 1}^{n} |x_i|
    \]

    Очевидно, что относительно такой нормы пространство будет полно -- из фундаментальности по $\|\cdot\|_e$ следует и покоординатная фундаментальность, а значит полнота следует из полноты $\R$.

    С другой стороны
    \[
        \|x\| \leq \sum_{i = 1}^{n} |x_i|\|e_i\| \leq \max_{i} \|e_i\| \sum_{i = 1}^{n} |x_i| = \max_{i} \|e_i\| \|x\|_{e} = C_1\|x\|_{e}.
    \]

    Пусть не существует $C_2$ такого, что
    \[
        \forall x \in X \hookrightarrow C_2 \|x\|_e \leq \|x\|,
    \]
    то есть $\forall C > 0 \ \exists x_C \neq 0 \in L\colon C \|x_C\|_e > \|x_C\|$.
    Рассмотрим тогда $C_n = \frac{1}{n}$ и $y_n = \frac{x_{C_n}}{\|x_{C_n}\|_e}$.
    И мы получаем, после подстановки, что $\|y_n\| < \frac{1}{n}$.

    Поскольку $|(y_n)_k| \leq 1$, то существует сходящаяся подпоследовательность $\{(y_{n_s})_k\}$ к некоторому $z_k\colon$
    \[
        (y_{n_s})_k \rightarrow z_k, s \rightarrow \infty.
    \]
    Выделяя эту подпоследовательность $n$ раз, мы получим сходимость к некоторому $z\colon \|z\|_e = 1$.

    Таким образом мы пришли к противоречию: с одной стороны, $y_n \xrightarrow{\|\cdot\|} 0$, а с другой -- $y_{n_s} \xrightarrow{\|\cdot\|_e} z$, где $z \neq 0$.
    Но из сходимости по $\|\cdot\|_e$-норме следует и сходимость по $\|\cdot\|$-норме, то есть у нас $y_{n_s} \xrightarrow{\|\cdot\|} z$ -- противоречие.

    Следовательно, фундаментальность по $\|\cdot\|$ и $\|\cdot\|_e$ эквивалентна, а следовательно $(L, \|\cdot\|)$ также является банаховым пространством.
\end{proof}
\begin{note}
    В частности, при доказательстве утверждения показано, что в конечномерных пространствах все нормы эквивалентны.
\end{note}
\begin{theorem}
    Пусть $(X, \|\cdot\|)$ -- нормированное пространство и $L \subset X$ -- его конечномерное подпространство.
    Тогда $L$ -- замкнуто.
\end{theorem}
\begin{proof}
    Поскольку $(L, \|\cdot\|\bigr|_{L})$ -- банахово пространство, то оно является и замкнутым.
    Покажем это.

    Пусть $x \in X$ -- точка прикосновения $L$ в $(X, \|\cdot\|)$.
    Тогда существует $\{x_n\} \subset L\colon x_n \rightarrow x$.
    Но $\{x_n\}$ -- фундаментальна, а, следовательно, сходится в $L$.
    В силу единственности предела -- $x \in L$.

\end{proof}
\section{Эквивалентность норм.}
\begin{definition}
    Пусть $X$ -- линейное пространство, и $\|\cdot\|_1$ с $\|\cdot\|_2$ -- две нормы на нём.
    Будем говорить, что $\|\cdot\|_1$ слабее $\|\cdot\|_2$, если $\exists C > 0 \forall x \in X \hookrightarrow \|x\|_1 \leq C\|x\|_2$.

    Будем говорить, что $\|\cdot\|_1$ эквивалентна $\|\cdot\|_2$, если $\|\cdot\|_1$ слабее $\|\cdot\|_2$ и $\|\cdot\|_2$ слабее $\|\cdot\|_1$.
\end{definition}
\begin{theorem}
    Если $X$ -- линейное пространство и $\|\cdot\|_1$ с $\|\cdot\|_2$ -- две нормы на нём, то следующие условия эквивалентны:
    \begin{enumerate}
        \item $\tau_1 \subset \tau_2$ (где $\tau_1$ и $\tau_2$ -- топологии, порождённые $\|\cdot\|_1$ и $\|\cdot\|_2$ соответственно).
        \item $\exists R > 0$ такое, что $O_R^2(0) \subset O_1^1(0)$.
        \item Норма $\|\cdot\|_1$ слабее нормы $\|\cdot\|_2$
    \end{enumerate}
\end{theorem}
\begin{proof}
    Пусть $\tau_1 \subset \tau_2$.
    Тогда $O_1^1(0) \in \tau_2$ и $0$ является внутренней точкой по $\tau_2$, а следовательно входит в $O_1^1(0)$ вместе с некоторым шаром $O_R^2(0)$.

    Пусть теперь $\exists R > 0$ такое, что $O_R^2(0) \subset O_1^1(0)$.
    В частности $B_{R/2}^2(0) \subset B_1^1(0)$.
    Тогда $\forall x \in X$ рассмотрим $z = \frac{Rx}{2\|x\|_2} \in B_{R/2}^2(0)$.
    Значит, оно лежит и в $B_1^1(0)$, то есть
    \[
        \|z\|_1 \leq 1.
    \]
    Подставим выражение для $z$ и получим:
    \[
        \|x\|_1 \leq \frac{2}{R}\|x\|_2.
    \]

    Пусть теперь $\|\cdot\|_1$ слабее $\|\cdot\|_2$.
    Тогда из сходимости последовательности по $\|\cdot\|_2$ следует сходимость последовательности и по $\|\cdot\|_1$, а значит если $A \in \phi_1$ (семейство замкнутых по $\|\cdot\|_1$), то если $x$ -- секвенциальная точка прикосновения по $\|\cdot\|_2$, то она является и секвенциальной точкой прикосновения по $\|\cdot\|_1$, а значит лежит в $A$ -- то есть $A \in \phi_2$ и, таким образом, всякое открытое по $\|\cdot\|_1$ открыто и по $\|\cdot\|_2$, что и показывает нужное включение.
\end{proof}

\begin{corollary}
    Пусть $X$ -- линейное пространство и $\|\cdot\|_1$ с $\|\cdot\|_2$ -- две нормы на нём.
    Тогда следующие утверждения эквивалентны:
    \begin{enumerate}
        \item $\tau_1 = \tau_2$.
        \item $\exists R_1, R_2 > 0\colon O_{R_1}^1(0) \subset O_1^2(0) \wedge O_{R_2}^2(0) \subset O_1^1(0)$.
        \item $\|\cdot\|_1$ и $\|\cdot\|_2$ эквивалентны.
    \end{enumerate}
\end{corollary}
\begin{proof}
    Очевидное следствие из предыдущего утверждения.
\end{proof}
\begin{proposition}
    Если $X$ -- линейное пространство и $\|\cdot\|_1$ с $\|\cdot\|_2$ -- две нормы на нём, то следующие условия эквивалентны:
    \begin{enumerate}
        \item $\tau_1 \subset \tau_2$ (где $\tau_1$ и $\tau_2$ -- топологии, порождённые $\|\cdot\|_1$ и $\|\cdot\|_2$ соответственно).
        \item Норма $\|\cdot\|_1$ слабее нормы $\|\cdot\|_2$.
        \item Из фундаментальности по $\|\cdot\|_2$ следует фундаментальность по $\|\cdot\|_1$.
        \item Из сходимости по $\|\cdot\|_2$ следует сходимость по $\|\cdot\|_1$.
    \end{enumerate}
\end{proposition}
\begin{proof}
    Эквивалентность $1$ и $2$ показана ранее.

    Импликация $2) \Rightarrow 3)$ -- очевидна.

    Пусть верно $3)$, но $2)$ -- неверно.
    Тогда, в частности, верно
    \[
        \forall k \in \N \ \exists x_k \in X\colon \|x\|_1 > k\|x\|_2.
    \]
    Рассмотрим $y_k = \frac{\sqrt{k} x_k}{\|x_k\|_1}$.
    Для последовательности $y_k$ верно, что
    \[
        \|y_k\|_1 = \sqrt{k}, \quad \|y_k\|_2 < \frac{1}{\sqrt{k}}.
    \]
    Таким образом $\{y_k\}$ -- фундаментальна по $\|\cdot\|_2$, а, следовательно, фундаментальна и по $\|\cdot\|_1$.
    Но всякая последовательность норм фундаментальной последовательности -- ограничена, что в этом случае не выполняется -- противоречие.

    Импликация $2) \Rightarrow 4)$ -- очевидна.
    Покажем теперь обратную импликацию.
    Используем аналогичную конструкцию.
    По $\|\cdot\|_2$ последовательность сходится к $0$, но при этом она не фундаментальна по $\|\cdot\|_1$ и, следовательно, не является сходящейся -- противоречие.
\end{proof}
\begin{example}[Пример последовательности, сходящейся к двум разным точкам по двум разным нормам.]
    Рассмотрим пространство $l^1$ с нормами $\|\cdot\|_*$ и $\|\cdot\|_{**}$, определяющимися как
    \[
        \|x\|_* = \sum_{n = 1}^{\infty} \frac{|x(n)|}{n}, \quad \|x\|_{**} = \left|\sum_{n = 1}^{\infty} x(n)\right| + \sum_{n = 2}^{\infty} \frac{|x(n)|}{n}.
    \]
    Рассмотрим последовательность $\{e_n\}_{n = 1}^{\infty}$ -- стандартные базисные векторы.
    Тогда по $\|\cdot\|_{*}\colon$
    \[
        \|e_n\|_{*} = \frac{1}{n} \rightarrow 0.
    \]
    А по $\|\cdot\|_{**}\colon$
    \[
        \|e_n - e_1\|_{**} = \frac{1}{n} \rightarrow 0.
    \]
\end{example}
\begin{note}
    Заметим, что для линейных пространств полнота по одной из эквивалентных норм равносильно полноте по другой.
\end{note}